\section{Orchestration}
\label{di:orchestration}
\subsection{Design}
A WinUI3 \texttt{MainWindow} starts as a single XAML file with a single code-behind translation unit. Although an MVVM or MVC architecture was considered, a 1-1 mapping of UI view and back-end library can be formed negating the requirement for a controller in this implementation. Use of a controller or view model would have increased the amount boilerplate for very little benefit. Therefore, the \texttt{MainWindow} class spreads its definition across several translation units:

\subsection{Implementation}
The decoupling of the \texttt{MainWindow} class is as follows

\begin{itemize}
    \item \texttt{MainWindow.Initialisation.cpp} owns \texttt{InitialiseTranscription()} and the launches the summarisation and RAG loaders.
    \item \texttt{MainWindow.Recording.cpp} owns the start/stop/cancel recording handlers and the processing-thread lambda that handles the audio recording to summarisation pipeline
    \item \texttt{MainWindow.Saving.cpp} owns transcript, summary and appraisal export through the WinRT \texttt{FileSavePicker}.
    \item \texttt{MainWindow.State.cpp} owns the \texttt{SetAppState()} state machine implementation.
    \item \texttt{MainWindow.Appraisal.cpp} and \texttt{MainWindow.DisplayAppraisals.cpp} own the JSON appraisal store and the corresponding UI list view.
    \item \texttt{MainWindow.Guidance.cpp} owns the file picker for adding PDFs, the indexer driver and the View Guidance click handler.
    \item \texttt{MainWindow.Settings.cpp} owns the persistent-settings flyout menu.
\end{itemize}


